% Generated draft: requires mathematical and source review.
\documentclass[11pt]{article}
\usepackage[margin=25mm]{geometry}
\usepackage{fontspec}
\setmainfont{Latin Modern Roman}
\usepackage{amsmath,amssymb,mathtools,unicode-math}
\setmathfont{Latin Modern Math}
\usepackage{enumitem,needspace,xurl,hyperref}
\setlength{\parindent}{0pt}
\setlength{\parskip}{5pt}
\setlength{\emergencystretch}{3em}
\newcommand{\N}{\mathbb{N}}
\newcommand{\Z}{\mathbb{Z}}
\newcommand{\Q}{\mathbb{Q}}
\newcommand{\R}{\mathbb{R}}
\newcommand{\C}{\mathbb{C}}
\newcommand{\F}{\mathbb{F}}
\newcommand{\A}{\mathbb{A}}
\newcommand{\PP}{\mathbb{P}}
\newcommand{\E}{\mathbb{E}}
\newcommand{\eps}{\varepsilon}
\newcommand{\dd}{\,\mathrm{d}}
\begin{document}
\section*{UnsolvedMath ID: 800257}
\textbf{Source problem code:} catalogue entry
\subsection{Definitions and mathematical statement}
Let \(\Delta=\{z\in\C:|z|<1\}\) be the open unit disc. The supplied problem asks: \emph{Classify proper holomorphic correspondences from \(\Delta\) to itself.}

A possible precise convention is to represent such a correspondence by a nonempty closed complex analytic subset \(\Gamma\subset\Delta\times\Delta\) of pure complex dimension one, with both coordinate projections
\[
\pi_1:\Gamma\to\Delta,\quad (z,w)\mapsto z,
\qquad
\pi_2:\Gamma\to\Delta,\quad (z,w)\mapsto w
\]
proper, surjective, and with finite fibers. Here a map is proper if the inverse image of every compact subset of its target is compact. The multivalued image of \(z\in\Delta\) is
\[
\Gamma(z)=\{w\in\Delta:(z,w)\in\Gamma\}.
\]
Under this convention, the requested task is to classify all such analytic subsets \(\Gamma\).

This is a conditional formalization: the supplied record does not define ``proper holomorphic correspondence.'' In particular, it does not specify the precise requirements on the coordinate projections, whether correspondences carry multiplicities, whether reducible or disconnected correspondences are included, or whether classification is literal or taken up to an equivalence relation. These choices remain unresolved; the convention above is not asserted to be the source's intended convention.
\subsection{Short English statement}
Classify proper holomorphic correspondences from the open unit disc to itself. A conditional interpretation uses one-dimensional analytic relations with both coordinate projections proper, surjective, and finite-to-one. The source does not specify the correspondence convention or the intended meaning of classification.
\subsection{Sources}
\begin{itemize}
\item Local UnsolvedMath registry, permanent ID \(800257\), source problem code \texttt{catalogue entry}, supplied complete source record in the AMR Open Problem Lists collection. The record attributes the problem to Buzzard, Fornaess, and Krantz, with proposed year 1995, and labels it \texttt{open}; this is supplied registry status, not an independently verified mathematical assessment.
\item Buzzard, Fornaess, and Krantz, \emph{The Several Complex Variables Problem List} (1995), Problem 56. Supplied source URL: \url{https://arxiv.org/abs/math/9509202}. The supplied provenance records access on 2026-10-16 and extraction by \texttt{arxiv\_latex}; these are source-record provenance claims and were not independently verified here.
\end{itemize}
\end{document}
